\subsection{The target in the rank-four construction at the rational level}
\label{sec:count-three-target}

Ordinary support has reached the target at the two chosen levels. Its Frobenius identities now allow the full-block exclusion and the rational conductor bounds needed to count the remaining representations.

\begin{proposition}[Rank-four bound at the original level]\label{prop:count-three-host-output}
For the data of Lemma~\ref{hyp:count-three-host}, with \(j\ge2\), and
using the four-link support of Proposition~\ref{prop:count-three-extra-vertices},
the characteristic-zero full blocks are excluded by the determinant test
of Proposition~\ref{prop:count-two-full-block-excluded}. Irreducible
targets have the rational flag, quadratic packet attachment and
level \(M_{\mathrm{host}}=\ell\,2^{n_{\rm III}}\mathcal Np_1\), with
\(n_{\rm III}\) defined in \eqref{eq:ordinary-dyadic-exponents}. Their normalized cohomology length satisfies

\[
 \log\ell\,s_\ell(\xi,j)\le(\ell\mathcal N j)^C
\]
 with the
constant independent of the fitted fields, exceptional places
and coefficient extension.
\end{proposition}

\begin{proof}
\proofstep{Classical control and full-block exclusion.}

After the four-link propagation the actual target kills
the sum of the two ordinary ideals. On the chosen CM
layer it has the target exponents \(j+1,j,1,0\), plus
the common link exponent. The common twist cancels in successive differences, giving
\[
 (j+1)-j=1,\qquad j-1\ge1,\qquad 1-0=1
\]
for \(j\ge2\). Subtracting the index from
the oppositely signed Hodge numbers gives a dominant
weight, exactly as in Proposition~\ref{prop:two-link-classical-control}.
The ordinary control argument there puts each evaluation
in the classical part of the same quotient. At the
exceptional prime of each link the residual matrix is
scalar, whereas the residual cyclotomic character has
nonidentity value. The rank-four tower supplies Frobenius elements with these two
residual properties at primes chosen separately from the seed prime. The Speh
argument of Proposition~\ref{prop:count-two-full-block-excluded} therefore
applies. The Frobenius modulus ratios at good primes are
strictly less than the residue norm.

For a full block on the
rational chart, first apply Proposition~\ref{prop:full-block-extension}
and run the propagation on the resulting nonsplit point.
The same ratio bound holds for the split point since the
characteristic polynomials agree. Adding the inertial exponents within each block gives its determinant exponent,
\[
 (j+1)+1=j+2,\qquad j+0=j.
\]
At a split good prime the ratio of determinant moduli is
\[
 p^{j+2}/p^j=p^2.
\]
For roots \(\lambda_i,\mu_i\) of the
respective blocks the two-link bound instead gives
\[
              \frac{|\lambda_1\lambda_2|}{|\mu_1\mu_2|}
                   =|\lambda_1/\mu_1|\,|\lambda_2/\mu_2|<p^2 .
\]
This is the required contradiction. The degree-zero and length comparisons of
Section~\ref{sec:rank-four-deformation} use this exclusion
on the chart for the rank-four construction.

\proofstep{Glue the two cuspidal attachments.}

At an irreducible target the four restricted residual
labels are distinct at the powered normalizer, so the
trace self-twist argument of
Proposition~\ref{prop:count-two-cuspidal-attachment} gives irreducibility
on \(G_L\). The two cuspidal transfers have equal polynomials
on the common good primes. Strong multiplicity one identifies
them. At each exceptional pair use the other level's local
comparison, and at \(\ell\) retain the comparison on the
marked side. Remove the common scalar twist to obtain the
cuspidal attachment of \(r|_{G_L}\).
\cite[Theorem 1.2 and Corollary 1.3]{newton-thorne-unitary}

\proofstep{The rational local ratio test.}
De Rhamness at \(\ell\) follows by descent from \(L\).
On the ordinary flag over \(\Q_\ell\), the four characters and their pairwise ratios are
\[
 \kappa_i=\nu_i\epsilon_\ell^{m_i},\qquad
 \frac{\kappa_i}{\kappa_h}=\frac{\nu_i}{\nu_h}\epsilon_\ell^{m_i-m_h},
 \qquad m=(j+1,j,1,0),
\]
with finite inertia for the \(\nu_i\). At an arithmetic Frobenius lift \(s\) with
\(\epsilon_\ell(s)=1\), the ordered residual labels
of the rational flag are
\[
                 (\bar\eta\bar\Psi,\bar\Psi,\bar\eta,1)(s).
\]
The adjacent characteristic ratios at \(s\) are
\[
 \bar\kappa_0/\bar\kappa_1=\bar\eta,\qquad
 \bar\kappa_1/\bar\kappa_2=\bar\Psi/\bar\eta,\qquad
 \bar\kappa_2/\bar\kappa_3=\bar\eta .
\]
The rational Frobenius test of Paragraph~\ref{par:dyadic-residual}
makes each of them nontrivial there.
This includes the middle pair when \(j=2\).
The Weil lines on the flag have parameters
\(\iota(\nu_i(g))\,\ell^{m_i\deg(g)}\) on the rational
Weil group.

Write \(\mu_i(g)=\iota(\nu_i(g))\ell^{m_i\deg(g)}\).
If two lines were related by the norm character, take the
order in which \(\mu_i(g)/\mu_h(g)=\ell^{\deg(g)}\).
Applying the coefficient-field \(\ell\)-valuation at a
Frobenius lift gives
\[
       m_i-m_h=1,\qquad
       \nu_i(g)/\nu_h(g)=1\quad\text{on the Weil group}.
\]
In the first computation the \(\nu_i\) are units. Weil-group
density then identifies the two finite-inertia factors. The residual values at \(s\) rule this out
for every pair with exponent difference one.

Indeed, on two such Weil lines of characters \(\mu_i,\mu_h\) the monodromy equation gives
\[
 (\mu_i(g)/\mu_h(g)-|g|)N_{ih}=0\qquad(g\text{ in the Weil group}),
\]
with the inverse norm in the opposite convention. The preceding exclusion in both orders forces every \(N_{ih}=0\).
Thus \(N=0\) rationally, and \(N|_{G_{L_v}}=0\) after restriction.
This use of the rational flag precedes restriction:
the quadratic root \(\bar\alpha\) is trivial at the
selected decomposition groups above \(\ell\) in \(L\),
where the rational separation need not hold.
\cite{colmez-fontaine}\label{cor:ordinary-count-iii-point-output-1}

At the selected places above \(\ell\) in \(L\), the
cuspidal transfer is tempered after undoing the common
algebraic normalization. As its parameter has zero
monodromy, the absolute values of all the ratios of
the Weil lines there are one. Scalar twisting cancels
in a ratio. Write \(f_{ih}\) for a ratio of two rational Weil lines.
For a rational Weil element \(g\), choose \(a>0\) with
\(g^a\) in the finite-index local restriction from \(L\).
Then
\[
                         |f_{ih}(g)|^a=|f_{ih}(g^a)|=1 .
\]
This proves the modulus equality on the rational group.
At the rational lift \(s\), \(\epsilon_\ell(s)=1\).
If the ratio of the finite-inertia factors were torsion at that
element, its complex absolute value would be one. This would give
\[
       |\mu_i(s)/\mu_h(s)|=\ell^{(m_i-m_h)\deg(s)}\ne1,
\]
contrary to the modulus equality.
Thus all these ratios have infinite order, and the same
is true on a positive power in any prescribed finite
extension's decomposition group.

\proofstep{Rigidity, descent and the valuation.}

Classicality supplies Frobenius data on the fitted field. Bounding the original character component requires counting the descended representations and their integral congruences.

The normalizer in \(G_L\) has trivial mod-\(\ell\)
cyclotomic character and retains four distinct labels.
Suppose a sign \(\chi\) through cyclotomic reduction
twists \(\bar\kappa_i\) to \(\bar\kappa_{\sigma(i)}\).
Writing \(g_0\) for this normalizer gives
\[
 \bar\kappa_{\sigma(i)}(g_0)=\chi(g_0)\bar\kappa_i(g_0)
 =\bar\kappa_i(g_0)\quad\Longrightarrow\quad\sigma(i)=i .
\]
The identity \(\chi\bar\kappa_i=\bar\kappa_i\) now holds as characters, so cancellation gives \(\chi=1\). The simple-spectrum
element just constructed and the argument of
Proposition~\ref{prop:marked-invariant-rigidity} give irreducibility on
\(G_{L(\mu_{\ell^\infty})}\).
The same adjoint vanishing theorem and local tangent
argument apply to the ordinary cuspidal attachment.
In particular the infinite-order ratios apply on the
open local subgroups needed for the weight-gap-one
extensions. The rational tangent class then vanishes
in the normalized chart.

\proofstep{Return to the original level.}

For the original field level take
\[
 M_{\mathrm{host}}=\ell\,2^{n_{\rm III}}\mathcal Np_1.
\]

Away from \(\ell\) the rational chart factors inertia
through the residual character quotient because
\(S_b=\varnothing\). It gives finite abelian inertia
with conductors dividing \(M_{\mathrm{host}}\).
At \(\ell\) the finite parts of the ordered rational
inertia characters have conductors dividing the same
level, and the rational ratio test supplies the zero
monodromy condition for this order of weights.
For quadratic descent choose the ramified prime of
\(K\) in the fitted tower outside the fixed packet
discriminants as in Proposition~\ref{prop:count-two-quadratic-packets}.
This is allowed in the simultaneous splitting
conditions of Section~\ref{sec:count-three-tower}.

At inertia at this ramified prime the original representation acts trivially and the quadratic sign acts nontrivially. The incompatible characteristic polynomials
\[
 (X-1)^4\quad\text{and}\quad (X+1)^4
\]
give the self-twist exclusion for each admissible quadratic restriction as in Proposition~\ref{prop:count-two-quadratic-packets}. Cyclic descent
and unramified comparison give the rationally
normalized cuspidal attachment on each member of the
packet list for which the restriction is irreducible.
At \(w\mid p\) of a member \(K'\) of the list, a
restricted unit character is the rational one composed
with \(N_{K'_w/\Q_p}\). For \(u\in1+M_{\mathrm{host}}\mathcal O_{K'_w}\), use the
matrix of multiplication in an integral local basis:
\[
 N_{K'_w/\Q_p}(u)=\det(I+M_{\mathrm{host}}A)
                         \in1+M_{\mathrm{host}}\Z_p .
\]
This gives triviality of the inducing unit characters on
the required congruence group. The weights on that side
are \(j+1,j,1,0\).
\cite{clozel-rajan,blggt,blggt-lg,thorne-adjoint}\label{cor:ordinary-count-iii-point-output-2}

\proofstep{Apply the pointwise valuation bound.}

The indexed trace product of the pointwise bound
therefore applies to the rank-four case with this \(M_{\mathrm{host}}\).
Choose \(p\geq5\) prime to \(2\ell\mathcal Np_1\),
with \(p\leq6\ell\mathcal Np_1+1\).
The reducible target and all other points have the
same trace comparison as in the even-order case, with the
\(2^{n_{\rm III}}\)-th roots included.
Each irreducible point maps to a distinct record in the fixed
packet product \(D_p\), with the two possible rational descents
retained as separate labels.

The cardinality bound of
Proposition~\ref{prop:marked-rational-packets} is polynomial in
\(M_{\mathrm{host}}(j+1)\).
The norm, integrality and length arguments then give
\[
      \log\ell\,s_\ell(\xi,j)\leq(\ell\mathcal N j)^C .
\]
Here \(p_1\) is fixed by the primitive target, and
\(2^{n_{\rm III}}\) and the order of the root group
in the trace are bounded by fixed powers of
\(\ell\mathcal N\). Thus \(C\) is independent of the
fitted fields, the weight of the Eisenstein seed, its
prime \(w\), both exceptional places and the degree of
the coefficient extension. This gives the rank-four positive-component valuation estimate with the same host and
parameters as the common-window calculation.
\end{proof}
